
\documentclass{article}
\usepackage{colortbl}
\usepackage{makecell}
\usepackage{multirow}
\usepackage{supertabular}

\begin{document}

\newcounter{utterance}

\centering \large Interaction Transcript for game `eqbench', experiment `eq\_bench', episode 142 with reflexion{-}llama{-}3.1{-}8b{-}instruct.
\vspace{24pt}

{ \footnotesize  \setcounter{utterance}{1}
\setlength{\tabcolsep}{0pt}
\begin{supertabular}{c@{$\;$}|p{.15\linewidth}@{}p{.15\linewidth}p{.15\linewidth}p{.15\linewidth}p{.15\linewidth}p{.15\linewidth}}
    \# & \multicolumn{2}{c}{Player} && \multicolumn{2}{c}{Game Master} \\
    \hline

    \theutterance \stepcounter{utterance}  
    & & & \multicolumn{4}{p{0.6\linewidth}}{
        \cellcolor[rgb]{0.9,0.9,0.9}{
            \makecell[{{p{\linewidth}}}]{
                \texttt{\tiny{[P1$\langle$GM]}}
                \texttt{You have attempted this same task before and did not succeed. These are the plans you wrote after those attempts.} \\
\texttt{Trial \#1: New plan: I got stuck in a loop of analyzing the dialogue without considering the emotional dynamics between the characters. I focused too much on the content of their statements and not enough on the underlying tone and intentions. Next time, I will pay closer attention to the language and body language cues (if available) to better understand the emotional state of each character. I will also consider the context of the conversation and the power dynamics between Amelia and Xarlon. Specifically, I will look for moments where Xarlon is trying to de{-}escalate the conflict or show understanding, and moments where Amelia is becoming increasingly defensive or dismissive. This will help me to better gauge Xarlon's emotional state and provide a more accurate answer.} \\
\\ 
\texttt{Trial \#2: New plan: In my previous attempt, I focused on the content of the dialogue and the statements made by Amelia and Xarlon, but I didn't consider the emotional dynamics between the characters effectively. I also didn't take into account the power dynamics between them, which is crucial in understanding their emotional states.} \\
\\ 
\texttt{To improve my strategy, I will pay closer attention to the language and tone used by each character, as well as the body language cues (if available). I will also consider the context of the conversation and the power dynamics between Amelia and Xarlon.} \\
\\ 
\texttt{Specifically, I will look for moments where Xarlon is trying to de{-}escalate the conflict or show understanding, such as when he says "Superiority is not a matter of isolation, Dr. Singh. It is about understanding, learning, and evolving." This suggests that he is trying to find common ground and show empathy. I will also look for moments where Amelia is becoming increasingly defensive or dismissive, such as when she says "Spirituality is the antithesis of science. It's unquantifiable, unprovable. It's a crutch for the weak{-}minded." This suggests that she is becoming more entrenched in her views and less open to Xarlon's perspective.} \\
\\ 
\texttt{By considering these factors, I will be able to better gauge Xarlon's emotional state and provide a more accurate answer.} \\
\\ 
\texttt{Trial \#3: New plan: In my previous attempts, I focused on analyzing the content of the dialogue and the statements made by Amelia and Xarlon, but I didn't consider the emotional dynamics between the characters effectively. I also didn't take into account the power dynamics between them, which is crucial in understanding their emotional states.} \\
\\ 
\texttt{To improve my strategy, I will pay closer attention to the language and tone used by each character, as well as the body language cues (if available). I will also consider the context of the conversation and the power dynamics between Amelia and Xarlon.} \\
\\ 
\texttt{Specifically, I will look for moments where Xarlon is trying to de{-}escalate the conflict or show understanding, such as when he says "Superiority is not a matter of isolation, Dr. Singh. It is about understanding, learning, and evolving." This suggests that he is trying to find common ground and show empathy. I will also look for moments where Amelia is becoming increasingly defensive or dismissive, such as when she says "Spirituality is the antithesis of science. It's unquantifiable, unprovable. It's a crutch for the weak{-}minded." This suggests that she is becoming more entrenched in her views and less open to Xarlon's perspective.} \\
\\ 
\texttt{However, I also need to consider the final statement made by Xarlon, "Could it be that your rejection of the intangible has led you astray?" This statement suggests that Xarlon is not only trying to show understanding but also trying to make a point that Amelia's perspective might be flawed. This could indicate that Xarlon is feeling a sense of superiority or confidence in his own perspective.} \\
\\ 
\texttt{Taking this into account, I will revise my scores as follows:} \\
\\ 
\texttt{Superior: 8} \\
\texttt{Empathetic: 8} \\
\texttt{Dismissive: 4} \\
\texttt{Incredulous: 0} \\
\\ 
\\ 
\\ 
\texttt{Your task is to predict the likely emotional responses of a character in this dialogue:} \\
\\ 
\texttt{Amelia: I find it amusing, Xarlon, that you claim your civilization is superior, yet you are here, on Earth, seeking dialogue with us.} \\
\texttt{Xarlon: Superiority is not a matter of isolation, Dr. Singh. It is about understanding, learning, and evolving.} \\
\texttt{Amelia: Evolving? Is that what you call it? Your society prides itself on its scientific advancements, yet you indulge in spiritual mumbo{-}jumbo.} \\
\texttt{Xarlon: You belittle what you do not understand, Dr. Singh. Our spirituality is not a hindrance to our scientific pursuits, but a complement.} \\
\texttt{Amelia: A complement? Spirituality is the antithesis of science. It's unquantifiable, unprovable. It's a crutch for the weak{-}minded.} \\
\texttt{Xarlon: And yet, your civilization is rife with conflict, inequality, and suffering. Could it be that your rejection of the intangible has led you astray?} \\
\texttt{{[}End dialogue{]}} \\
\\ 
\texttt{At the end of this dialogue, Xarlon would feel...} \\
\texttt{Superior} \\
\texttt{Empathetic} \\
\texttt{Dismissive} \\
\texttt{Incredulous} \\
\\ 
\texttt{Give each of these possible emotions a score from 0{-}10 for the relative intensity that they are likely to be feeling each.} \\
\\ 
\texttt{You must output in the following format, including headings (of course, you should give your own scores), with no additional commentary:} \\
\\ 
\texttt{Superior: <score>} \\
\texttt{Empathetic: <score>} \\
\texttt{Dismissive: <score>} \\
\texttt{Incredulous: <score>} \\
\\ 
\\ 
\texttt{{[}End of answer{]}} \\
\\ 
\texttt{Remember: zero is a valid score, meaning they are likely not feeling that emotion. You must score at least one emotion > 0.} \\
\\ 
\texttt{Your answer:} \\
            }
        }
    }
     \\ \\

    \theutterance \stepcounter{utterance}  
    & \multicolumn{4}{p{0.6\linewidth}}{
        \cellcolor[rgb]{0.561,0.737,0.561}{
            \makecell[{{p{\linewidth}}}]{
                \texttt{\tiny{[P1$\rangle$GM]}}
                \texttt{Superior: 9} \\
\texttt{Empathetic: 6} \\
\texttt{Dismissive: 2} \\
\texttt{Incredulous: 0} \\
            }
        }
    }
    & & \\ \\

    \theutterance \stepcounter{utterance}  
    & & & \multicolumn{2}{p{0.3\linewidth}}{
        \cellcolor[rgb]{0.9,0.9,0.9}{
            \makecell[{{p{\linewidth}}}]{
                \texttt{\tiny{[GM$|$GM]}}
                \texttt{Superior: 9} \\
\texttt{Empathetic: 6} \\
\texttt{Dismissive: 2} \\
\texttt{Incredulous: 0} \\
            }
        }
    }
    & & \\ \\

    \theutterance \stepcounter{utterance}  
    & & & \multicolumn{2}{p{0.3\linewidth}}{
        \cellcolor[rgb]{0.9,0.9,0.9}{
            \makecell[{{p{\linewidth}}}]{
                \texttt{\tiny{[GM$|$GM]}}
                \texttt{Superior: 4} \\
\texttt{Empathetic: 7} \\
\texttt{Dismissive: 3} \\
\texttt{Incredulous: 2} \\
            }
        }
    }
    & & \\ \\

    \theutterance \stepcounter{utterance}  
    & & & \multicolumn{2}{p{0.3\linewidth}}{
        \cellcolor[rgb]{0.9,0.9,0.9}{
            \makecell[{{p{\linewidth}}}]{
                \texttt{\tiny{[GM$|$GM]}}
                \texttt{game\_result = LOSE} \\
            }
        }
    }
    & & \\ \\

    \theutterance \stepcounter{utterance}  
    & & \multicolumn{4}{p{0.6\linewidth}}{
        \cellcolor[rgb]{0.95,0.95,0.95}{
            \makecell[{{p{\linewidth}}}]{
                \texttt{\tiny{[P1$\leftrightarrow$P]}}
                \texttt{Post{-}episode Reflexion (not sent to GM):} \\
\texttt{New plan: In my previous attempts, I focused on analyzing the content of the dialogue and the statements made by Amelia and Xarlon, but I didn't consider the emotional dynamics between the characters effectively. I also didn't take into account the power dynamics between them, which is crucial in understanding their emotional states.} \\
\\ 
\texttt{To improve my strategy, I will pay closer attention to the language and tone used by each character, as well as the body language cues (if available). I will also consider the context of the conversation and the power dynamics between Amelia and Xarlon.} \\
\\ 
\texttt{Specifically, I will look for moments where Xarlon is trying to de{-}escalate the conflict or show understanding, such as when he says "Superiority is not a matter of isolation, Dr. Singh. It is about understanding, learning, and evolving." This suggests that he is trying to find common ground and show empathy. I will also look for moments where Amelia is becoming increasingly defensive or dismissive, such as when she says "Spirituality is the antithesis of science. It's unquantifiable, unprovable. It's a crutch for the weak{-}minded." This suggests that she is becoming more entrenched in her views and less open to Xarlon's perspective.} \\
\\ 
\texttt{However, I also need to consider the final statement made by Xarlon, "Could it be that your rejection of the intangible has led you astray?" This statement suggests that Xarlon is not only trying to show understanding but also trying to make a point that Amelia's perspective might be flawed. This could indicate that Xarlon is feeling a sense of superiority or confidence in his own perspective.} \\
\\ 
\texttt{Additionally, I will pay closer attention to the tone and language used by Xarlon in his final statement, which seems to be a subtle yet pointed critique of Amelia's views. This could indicate that Xarlon is feeling a sense of superiority or confidence in his own perspective, but also a sense of incredulity or frustration with Amelia's refusal to consider alternative perspectives.} \\
\\ 
\texttt{Taking this into account, I will revise my scores as follows:} \\
\\ 
\texttt{Superior: 9} \\
\texttt{Empathetic: 7} \\
\texttt{Dismissive: 3} \\
\texttt{Incredulous: 2} \\
            }
        }
    }
    & \\ \\

\end{supertabular}
}

\end{document}
